\documentclass[12pt]{article}
\usepackage[utf8]{inputenc}
\usepackage[T1]{fontenc}
\usepackage{hyperref}
\title{Central Pomo Language Resources for K-12 Curriculum Development}
\author{Compiled by Claude Sonnet 5.5 (claude-sonnet-5-5) for Andr\'as Kornai}
\date{}
\begin{document}
\maketitle
\section{Introduction}
Central Pomo (ISO 639-3 poo, Glottocode cent2138) is a Pomoan language of northern California, classed as dormant: no first-language speakers remain. Searches found very little material specific to Central Pomo; most Pomo-language resources concern other Pomoan languages (Southern, Eastern, Southeastern, Kashaya) and must not be mistaken for Central Pomo materials. This report is therefore short.
\section{Primary Educational Resources}
No Central Pomo primer, textbook or school programme was found in search results.
\section{Linguistic and Reference Materials}
\begin{itemize}
\item Glottolog entry for Central Pomo (classification and status): \url{https://glottolog.org/resource/languoid/id/cent2138}
\item Endangered Languages Project page for Central Pomo: \url{https://endangeredlanguages.com/node/37871}
\item Wikipedia overview of Central Pomo (mentions a transcription project on J.P. Harrington's Central Pomo materials): \url{https://wikipedia.com/wiki/ISO_639:poo}
\item UBC RelLex introduction to the Central Pomo language: \url{https://wiki.ubc.ca/Documentation:RelLex/Central_Pomo_Language_-_Intro}
\end{itemize}
\section{Religious and Community Resources}
\begin{itemize}
\item California Indian Museum and Cultural Center, Pomo Language Status Report (needs assessment for Pomo language preservation; covers the Pomo languages generally): \url{https://cimcc.org/wp-content/uploads/2012/05/Pomo-Language-Status-Report.pdf}
\end{itemize}
\section{Digital and Audio Resources}
No Central Pomo-specific audio or digital dictionary was found. Other Pomoan languages have such tools, but they are not Central Pomo and are not listed.
\section{Curriculum Development Context}
With no living first-language speakers and little published material, a Central Pomo curriculum would depend on archival sources (such as the Harrington notes) and on cooperation with Pomo community organisations. Contact the community bodies above before any development work.
\end{document}